\chapter{תוכנית אימות ובדיקות}\label{ch:verif}

כל דרישה בעדיפות חובה נבדקת לפחות פעם אחת, לפי השיטה שנקבעה לה בפרק~\ref{ch:req}. מטריצת העקיבות המלאה נמצאת בנספח~\ref{app:trace}.

\section{מאגר הבדיקה \en{VL-Wild}}\label{sec:bench}
המאגר הוא אבן הבוחן של הפרויקט, ולכן הוא מוגדר כבר בדוח המכין ולא ניתן לשינוי בלי בקשת שינוי:
\begin{itemize}
\item \textbf{גודל:} \num{20,000} תמונות, \num{10,000} אמיתיות ו-\num{10,000} סינתטיות.
\item \textbf{מקור האמיתיות:} צילומי חדשות מתוך \num{12} החודשים האחרונים, מהסוכנויות שמהן נרכש מאגר האימון (\req{A-02}), אבל תמונות שלא נכללו באימון.
\item \textbf{מקור הסינתטיות:} \num{50\%} תמונות שאספו בודקי עובדות ״מהשטח״ (כלומר הופצו בפועל ברשת), ו-\num{50\%} שנוצרו בשלושה מחוללים שהוצאו מהאימון.
\item \textbf{עיבוד:} כל התמונות עוברות את שרשרת ההפצה שבפרק~\ref{ch:method}.
\item \textbf{בעלות:} המאגר נשמר אצל צד שלישי (המוסד האקדמי) ונחשף לצוות רק ביום הבדיקה, כדי למנוע ״התאמת יתר״ אליו.
\end{itemize}

\section{תרחישי בדיקה מרכזיים}
\small

לכל תרחיש בטבלה~\ref{tab:tests} נקבע קריטריון מעבר מראש, כדי שלא נחליט רק אחרי הבדיקה מה נחשב הצלחה.

\begin{longtable}{@{}l p{5.6cm} p{4.7cm} l@{}}
\caption{תרחישי בדיקה עיקריים וקריטריון מעבר}\label{tab:tests}\\
\toprule
\textbf{מזהה} & \textbf{תרחיש} & \textbf{קריטריון מעבר} & \textbf{דרישה} \\
\midrule
\endfirsthead
\toprule
\textbf{מזהה} & \textbf{תרחיש} & \textbf{קריטריון מעבר} & \textbf{דרישה} \\
\midrule
\endhead
\bottomrule
\endfoot
\req{T-01} & הרצת המערכת על כל \en{VL-Wild} & גילוי \num{90\%} לפחות, התרעות שווא \num{2\%} לכל היותר & \req{PR-04} \\
\req{T-02} & אותה בדיקה אחרי דחיסה באיכות \num{70} והקטנה ל-\num{50\%} & ירידה של עד \num{5} נקודות בשיעור הגילוי & \req{PR-05} \\
\req{T-03} & \num{500} תמונות עם תעודת \en{C2PA} תקפה, \num{100} פגומות, \num{50} שבוטלו & \num{100\%} סיווג נכון של התעודה & \req{FR-04}, \req{FR-05} \\
\req{T-04} & מבחן כיול: \num{10} קבוצות ציון & הפרש ממוצע בין ציון לשיעור הזיוף $\leq$ \num{5} נק׳ & \req{FR-08} \\
\req{T-05} & ניסיון לפרסם תמונה ״לא ודאית״ בהרשאת עורך רגיל & הפרסום נחסם ונרשם ביומן & \req{FR-14} \\
\req{T-06} & ניסיון לשנות רשומת יומן ישירות במסד הנתונים & שינוי מתגלה בבדיקת שלמות & \req{FR-17} \\
\req{T-07} & עומס: \num{11,000} תמונות בשעה במשך \num{4} שעות & אחוזון \num{95} $\leq$ \num{180} שנ׳; אפס אובדן & \req{PR-01}, \req{PR-02} \\
\req{T-08} & הדגמת לוח בקרה לעורך ראשי & היסטוגרמה, \en{ROC} וגרף עמודות כמוגדר & \req{FR-19} \\
\req{T-09} & זמן החלטה בפיילוט, \num{4} שבועות & חציון $\leq$ \num{15} דק׳ & \req{G-2} \\
\req{T-10} & ביקורת מדגמית של \num{2,000} תמונות לפני ואחרי & ירידה $\geq$ \num{80\%} בתמונות שקריות שפורסמו & \req{G-3} \\
\req{T-11} & סריקת תוויות בכל עמודי האתרים בפיילוט & \num{100\%} תמונות מתויגות & \req{FR-15}, \req{G-4} \\
\req{T-12} & בדיקת חדירה בידי חברה חיצונית & אפס ממצאים קריטיים פתוחים & \req{SR-01}, \req{SR-02} \\
\end{longtable}
\normalsize

\section{עקומת היעד}
איור~\ref{fig:roc} מגדיר את היעד בצורה גרפית: עקומת \en{ROC} של המערכת חייבת לעבור דרך האזור המסומן או מעליו. העקומה המקווקוות היא דוגמה לגלאי בודד ״מהשטח״, קרוב לקו הניחוש האקראי, כפי שתיאר~\cite{ren2025}. עקומת \en{ROC} זו היא גם אחד התוצרים הגרפיים המחייבים בלוח הבקרה (\req{FR-19}).

\begin{figure}[!htb]
\begin{ltrdiagram}
% roc: TikZ / pgfplots source, built inside ltrdiagram (LTR)
\begin{tikzpicture}
\begin{axis}[width=8.4cm,height=7cm,xmin=0,xmax=0.3,ymin=0.5,ymax=1,
  xlabel={\en{FPR}}, ylabel={\en{TPR}},
  ticklabel style={font=\sffamily\scriptsize},
  xtick={0.02,0.1,0.2,0.3}, xticklabels={2\%,10\%,20\%,30\%},
  ytick={0.5,0.6,0.7,0.8,0.9,1}, yticklabels={50\%,60\%,70\%,80\%,90\%,100\%},
  legend style={font=\sffamily\scriptsize,at={(0.97,0.03)},anchor=south east},
  grid=both, grid style={gray!20}]
\fill[okc!20] (axis cs:0,0.9) rectangle (axis cs:0.02,1);
\addplot[very thick,acc2,smooth] coordinates {(0,0.5)(0.005,0.8)(0.01,0.87)(0.02,0.915)(0.05,0.95)(0.1,0.975)(0.2,0.99)(0.3,0.995)};
\addplot[thick,dashed,gray,smooth] coordinates {(0,0.5)(0.05,0.56)(0.1,0.61)(0.2,0.69)(0.3,0.76)};
\addplot[thin,black!40] coordinates {(0,0)(0.3,0.3)};
\legend{\he{יעד \en{VeriLens}, משוקלל},\he{גלאי בודד בתנאי שטח},}
\draw[okc!60!black] (axis cs:0.02,0.95)--(axis cs:0.06,0.88) node[anchor=north west,font=\sffamily\scriptsize]{\he{אזור היעד}};
\end{axis}
\end{tikzpicture}

\end{ltrdiagram}
\caption{עקומת \en{ROC}: המערכת חייבת לעבור באזור הירוק, הרחק מגלאי בודד}\label{fig:roc}
\end{figure}

\section{קבלת המערכת}
הקבלה נעשית בשני שלבים. \textbf{קבלה טכנית} בחודש \num{12}: הרצת \req{T-01} עד \req{T-08} בנוכחות המרצה או נציג/ה מטעמו. \textbf{קבלה מבצעית} בחודש \num{18}: \req{T-09} עד \req{T-11} על נתוני הפיילוט. כשל בבדיקה מחייב תיקון והרצה חוזרת תוך \num{30} יום. כשל שני מוביל לדיון בבקשת שינוי (סעיף~\ref{sec:change}).
